\section{H. Trace and determinant in cote 1}
\label{sec:geometry-H}

\subsection{The trace formula and the genus-zero lemma}
The crossed-out sheets 79--80 and sheets 111--112 are mathematically useful;
being crossed out neither proves nor disproves their assertions. These
references in the prompt are physical PDF page numbers; direct inspection
identifies the public PDF footers as 78--79 and 110--111, respectively.
The lower-left pencil numbers are the archival numbering; for the first
pair they are 80--81. These three systems must not be conflated. Here all
Hilbert spaces are complex, eigenvalues are counted with \emph{algebraic}
multiplicity, and zero eigenvalues may be omitted. Write $\|T\|_1$ for the
trace norm of a trace-class operator $T$.

\begin{lemma}[Proved: coefficient and eigenvalue bounds]
\label{geo-H-bounds}
For $a_n(T)=\operatorname{Tr}(\bigwedge^nT)$ and singular values $s_j(T)$,
\[
 |a_n(T)|\leq\|\bigwedge^nT\|_1
   =\sum_{j_1<\cdots<j_n}s_{j_1}(T)\cdots s_{j_n}(T).
\]
Consequently $D_T(z)=\sum_{n\geq0}a_n(T)z^n$ is entire and, for every
$\epsilon>0$, $|D_T(z)|\leq C_\epsilon e^{\epsilon|z|}$. Its nonzero
eigenvalues satisfy $\sum_j|\lambda_j(T)|\leq\|T\|_1$.
\end{lemma}
\begin{proof}
A singular-value decomposition of $T$ induces one of $\bigwedge^nT$; its
singular values are the indicated products. For a trace-class operator $V$,
$|\operatorname{Tr}V|\leq\|V\|_1$, by the polar decomposition and the
Cauchy--Schwarz inequality for the positive functional
$W\mapsto\operatorname{Tr}(|V|^{1/2}W|V|^{1/2})$. Thus
\[
 |D_T(z)|\leq h(|z|),\qquad h(r)=\prod_j(1+rs_j(T)).
\]
This product converges uniformly on compact sets because $\sum_js_j(T)$ is
finite. Choose $J$ with $\sum_{j>J}s_j<\epsilon/2$. Its tail is at most
$e^{\epsilon r/2}$, and the finite head is a polynomial bounded by
$C_\epsilon e^{\epsilon r/2}$. This proves the entire-function assertion.

For any finite collection of nonzero eigenvalues with multiplicity, compact
operator spectral theory provides a finite-dimensional invariant space $V$
spanned by the relevant generalized eigenvectors. If a multiplicity is only
partly selected, take initial segments of its Jordan chains. Let $P$ be the
orthogonal projection onto $V$. The compression $PT|_V$ has the selected
eigenvalues. A unitary Schur triangularization and a diagonal unitary matrix
of appropriate phases give
\[
 \sum_{\lambda\in V}|\lambda|
 \leq\|PT|_V\|_1\leq\|T\|_1.
\]
The first inequality follows because the trace of the phase-adjusted triangular
matrix is the sum of the absolute values of its diagonal entries; the second
is the trace-ideal inequality for multiplication by contractions. Taking the
supremum over finite collections proves summability. The weaker bound
$4\|T\|_1$ also follows from the manuscript's separate positive and negative
real and imaginary parts. No infinite-dimensional trace/eigenvalue identity
has been used in either argument.
\end{proof}

\begin{lemma}[Proved: the required genus-zero lemma]
\label{geo-H-genus}
Suppose $f$ is a nonzero entire function with $f(0)\ne0$, with
$|f(z)|\leq C_\epsilon e^{\epsilon|z|}$ for every $\epsilon>0$. If its
nonzero zeros $z_j$, with multiplicity, satisfy $\sum_j|z_j|^{-1}<\infty$,
then
\[
 f(z)=f(0)\prod_j(1-z/z_j).
\]
\end{lemma}
\begin{proof}
Hadamard factorization for order at most one gives
$f(z)=f(0)e^{az}g(z)$, where $g(z)=\prod_j(1-z/z_j)$.
Indeed the genus-one elementary factors can be replaced by genus-zero factors
because $\sum_j1/z_j$ converges absolutely; their linear exponential factors
are absorbed into $a$. The version used is Theorem 3.16 of
McMullen~\cite{McMullen2025}, an opened source. The finite-head/tail argument
of Lemma~\ref{geo-H-bounds} gives $|g(z)|\leq C'_\epsilon e^{\epsilon|z|}$.

If $a\ne0$, replace $z$ by $(\overline a/|a|)z$. Thus assume $a>0$.
For $|\arg z|\leq\pi/3$, choose $\epsilon=a/4$ in the bound for $f$ to get
$|g(z)|\leq C e^{-c|z|}$, with $c=a/4>0$. Globally use the bound
$|g(z)|\leq C'e^{c|z|/3}$. Put $\omega=e^{2\pi i/3}$. For every nonzero
$z$, at least one of $z,\omega z,\omega^2z$ has argument between
$-\pi/3$ and $\pi/3$. Therefore
\[
 |g(z)g(\omega z)g(\omega^2z)|\leq C(C')^2e^{-c|z|/3}.
\]
The left-hand product is an entire bounded function. Liouville's theorem and
its decay imply that it is identically zero, contradicting its value $1$ at
zero. Hence $a=0$.
\end{proof}

\begin{theorem}[Proved from the Fredholm zero theorem: Lidskii's formula]
\label{geo-H-lidskii}
Every trace-class operator $T$ on a complex Hilbert space satisfies
\[
 \det(1+zT)=\prod_j(1+z\lambda_j(T)),\qquad
 \operatorname{Tr}T=\sum_j\lambda_j(T),
\]
where the product converges locally uniformly and the series converges
absolutely.
\end{theorem}
\begin{proof}
For the determinant defined by its exterior-power coefficients, the Fredholm
zero theorem says that its zeros are precisely $-1/\lambda_j(T)$, with the
algebraic multiplicities of the corresponding nonzero eigenvalues. An exact
opened source is Grothendieck~\cite{Grothendieck1956Fredholm}, Chapter II,
\S4, Theorem 2 and its corollary, printed pp.~354--355, and Theorem 4,
printed pp.~359--361. His
nuclear-operator statement applies because a trace-class Hilbert operator
has its summable singular-value nuclear expansion. Its determinant is this
exterior-power determinant. These cited results supply only the zero locus
and its multiplicities, not the asserted product formula or trace identity.

Lemma~\ref{geo-H-bounds} verifies both hypotheses of
Lemma~\ref{geo-H-genus} for $D_T$; also $D_T(0)=1$. That lemma gives the
product formula. Uniform convergence of the product near zero and absolute
summability allow its derivative at zero to be taken term by term. The
coefficient of $z$ in $D_T$ is $a_1(T)=\operatorname{Tr}T$. Its product
coefficient is $\sum_j\lambda_j(T)$. Equating them proves the formula.
For a nonseparable Hilbert space, work on the separable reducing space
$\overline{\operatorname{ran}T+\operatorname{ran}T^*}$; on its orthogonal
complement $T$ is zero, which changes none of these quantities.
\end{proof}

On sheet~71, the proposed reduction to positive Hermitian operators is
only announced. Coefficient domination by the positive determinant does
not itself identify the entire function with its product over zeros:
one must still eliminate the possible exponential factor in Hadamard
factorization. Sheets~79--80 isolate that issue and sheets~111--112 supply
the missing genus-zero argument, proved above. Thus the gap on~71 is a
missing argument, not a refutation of the trace formula or of the
completed reduction.

The following are actual proof checks, not editorial preferences. In the
transcription of sheets 111--112, a sum in the displayed definition of $h$
must be a product for the proposed argument to apply; with a literal sum the
claimed entire majorant is not established. A coefficientwise majorization
gives $|f(z)|\leq h(|z|)$, not generally $|f(z)|\leq|h(z)|$. The rotation
that makes a nonreal exponent $a$ positive is $\overline a/|a|$, not
$a/|a|$. The angular interval must have $\pi/3\leq\theta<\pi/2$; the
larger interval reaching $\pi$ does not give positive $\cos\theta$.
Finally the third factor is $g(\omega^2z)$: the alternative reading
$g(\omega^3z)$ repeats the first direction and does not cover the plane.
The proof above treats the mathematically valid readings and identifies
exactly why each alternative fails. Ordinary eigenvectors alone do not count
algebraic multiplicities; generalized eigenvectors repair that point on
79--80.

\subsection{The infinite-trace definition is not the finite-trace definition}
Let $(\mathcal M,\tau)$ be a von Neumann algebra with a faithful normal
semifinite trace. Put $\mathcal J=\mathcal M\cap L^1(\tau)$ and
$\mathcal D=1+\mathcal J$. If $\tau(1)<\infty$, then
$\mathcal D=\mathcal M$. The following correction includes noninvertible
operators, but it uses the spectral trace rather than the manuscript's
unqualified integral of a decreasing rearrangement over an infinite interval.

\begin{example}[Refuted: the literal infinite rearrangement is multiplicative]
\label{geo-H-rearrangement}
Take $\mathcal M=B(\ell^2(\mathbf N))$ with its usual trace and let $P$ be
a rank-one projection. Set $A=1-\tfrac12P$ and $B=1+P$. These are positive
invertible operators in $1+\mathcal J$, and $AB=1$. The decreasing
rearrangement $\mu_t(A)$ is identically $1$: infinitely many singular values
are $1$, and the single value $1/2$ is invisible to this rearrangement.
For $B$, $\mu_t(B)=2$ on $0\leq t<1$ and is $1$ thereafter. Consequently
\[
 e^{\int_0^\infty\log\mu_t(A)\,dt}=1,\quad
 e^{\int_0^\infty\log\mu_t(B)\,dt}=2,\quad
 e^{\int_0^\infty\log\mu_t(AB)\,dt}=1.
\]
Thus the literal proposed formula is false already for invertible
finite-rank perturbations of $1$. The finite matrix block is verified by
\texttt{code/geometry\_determinant\_checks.py}; the infinite identity tail
explains why finite truncations of the \emph{rearrangement} do not compute the
same invariant. The spectral-trace determinant instead gives $1/2$, $2$,
and $1$, respectively.
\end{example}

\begin{theorem}[Proved: the corrected determinant, including noninvertibles]
\label{geo-H-determinant}
For $A\in\mathcal D$ define
\[
 \Delta_\tau(A)=\exp\bigl(\tau(\log|A|)\bigr),\qquad e^{-\infty}=0.
\]
The positive part of $\log|A|$ is trace integrable, so this is unambiguously
in $[0,\infty)$. For all $A,B\in\mathcal D$, including noninvertible ones,
\[
 \Delta_\tau(AB)=\Delta_\tau(A)\Delta_\tau(B),\qquad
 \Delta_\tau(A^*)=\Delta_\tau(A).
\]
It follows that, for $H=A^*A$,
\[
 \Delta_\tau(B^*HB)
 =\Delta_\tau(A^*A)\Delta_\tau(B^*B).
\]
When $\tau(1)<\infty$, these statements hold for every $A,B\in\mathcal M$.
\end{theorem}

Here is a proof with the limiting step supplied. The invertible case agrees
with Hochs--Kaad--Schemaitat~\cite{HKS2018}, Definition 7.4 and Propositions
7.6--7.7; that source was opened, but its statements concern invertibles.
They are not a citation for the noninvertible step below.

\begin{lemma}[Proved: trace-log regularization]
\label{geo-H-regularization}
For positive bounded $C$ with $C-1\in\mathcal J$, write
$\ell(C)=\tau(\log C)\in[-\infty,\infty)$. This functional is upper
semicontinuous for trace-norm convergence within this positive set. If
$D\geq0$, $D-1\in\mathcal J$, and $\ell(D)>-\infty$, then
\[
 \lim_{\epsilon\downarrow0}
 \ell\left(\frac{C+\epsilon D}{1+\epsilon}\right)=\ell(C).
\]
In particular this holds for $D=1$.
\end{lemma}
\begin{proof}
For $\delta>0$ let $f_\delta(s)=\log((s+\delta)/(1+\delta))$.
Then $f_\delta(1)=0$, $f_\delta(C)\in L^1(\tau)$, and the resolvent identity
and the integral formula for the logarithm give
\[
 |\tau(f_\delta(C')-f_\delta(C))|
 \leq\delta^{-1}\|C'-C\|_1.
\]
Indeed the trace norm of the difference of the two resolvents is at most
$\|C'-C\|_1(\delta+t)^{-2}$, whose integral for $t\geq0$ is
$\delta^{-1}\|C'-C\|_1$. On $0\leq s\leq1$, $\log s\leq f_\delta(s)$;
on $s\geq1$,
$\log s-f_\delta(s)\leq\delta(s-1)$. Therefore
\[
 \ell(C')\leq\tau(f_\delta(C'))+
                       \delta\|C'-1\|_1.
\]
For $C_n\to C$ in trace norm, take the upper limit and then let
$\delta\downarrow0$. On $s<1$ monotone convergence applies to the negative
parts of $f_\delta(s)$; on $s\geq1$ they are dominated by $s-1$.
Thus $\tau(f_\delta(C))\to\ell(C)$ and upper semicontinuity follows.

The matrices $(C+\epsilon D)/(1+\epsilon)$ converge to $C$ in trace norm,
giving the upper bound for the asserted limit. If $\ell(C)=-\infty$, this
already proves it. Otherwise $\log C$ and $\log D$ are trace integrable.
Operator concavity of the logarithm, followed by the trace, gives
\[
 \ell\left(\frac{C+\epsilon D}{1+\epsilon}\right)
 \geq\frac{\ell(C)+\epsilon\ell(D)}{1+\epsilon}.
\]
The same inequality when either operator is not bounded below follows by
first using $(C+\delta1)/(1+\delta)$ and
$(D+\delta1)/(1+\delta)$; their convex combination is normalized in the
same way. Each logarithm is then trace integrable, even when $\tau(1)$ is
infinite, and the spectral limit already proved for $f_\delta$ gives the
inequality as $\delta\downarrow0$. No scalar multiple of $\tau(1)$ is
subtracted. Operator concavity itself follows by integrating operator convexity
of $X\mapsto(X+t)^{-1}$ in the resolvent formula for $\log X$; thus it
applies inside any von Neumann algebra. This supplies the lower bound.
\end{proof}

\begin{proof}[Proof of Theorem~\ref{geo-H-determinant}]
First $A^*A-1\in\mathcal J$ and
$|A|-1=(A^*A-1)(|A|+1)^{-1}\in\mathcal J$.
The scalar inequality $\log^+s\leq|s-1|$ proves the asserted integrability.
Products remain in $\mathcal D$.

\emph{Adjoints.} The polar decomposition identifies the spectral measures
of $A^*A$ and $AA^*$ away from zero by a partial isometry, which preserves
$\tau$. Their kernel projections have finite trace because each is dominated
by $|A^*A-1|$ or $|AA^*-1|$, respectively. Writing $A=1+K$ gives
\[
 \tau(A^*A-1)=\tau(K+K^*+K^*K)
 =\tau(K+K^*+KK^*)=\tau(AA^*-1).
\]
The integrals of $s-1$ over the nonzero spectral parts agree and are
absolutely convergent. The displayed equality therefore forces the kernel
traces to agree as well. Integrating $\log s$, with its allowed infinite
negative part, proves $\ell(A^*A)=\ell(AA^*)$. Hence
$\Delta_\tau(A)=\Delta_\tau(A^*)$.

\emph{Invertibles.} If $U$ is invertible in $\mathcal D$, both
$\log|U|$ and a bounded self-adjoint argument $Q$ of its polar unitary lie
in $\mathcal J$. For the argument use
$|\theta|\leq(\pi/2)|e^{i\theta}-1|$ on $[-\pi,\pi]$; for
$\log|U|$ use the bounded-below spectrum of $|U|$. Thus
$U(t)=e^{itQ}e^{t\log|U|}$ joins $1$ to $U$ through invertibles, smoothly
in the norm $\|\cdot\|+\|\cdot\|_1$. Differentiating the resolvent formula
and using the trace property gives, for any such path $V(t)$,
\[
 \frac{d}{dt}\log\Delta_\tau(V(t))
  =\operatorname{Re}\tau(V(t)^{-1}V'(t)).
\]
For detail, differentiation of $P(t)=V(t)^*V(t)$ gives
$\frac d{dt}\tau\log P=\tau(P^{-1}P')$; cyclicity reduces half of this
to the displayed real part. Resolvent integrals converge in trace norm on
compact path intervals because $P(t)$ has a uniform positive lower bound.
For fixed invertible $B$, the derivatives of
$\log\Delta_\tau(U(t)B)$ and $\log\Delta_\tau(U(t))$ agree by cyclicity.
Evaluation at $t=0,1$ proves multiplicativity for invertibles.

\emph{One invertible factor.} For arbitrary $A\in\mathcal D$ put
$A_\epsilon=((A^*A+\epsilon1)/(1+\epsilon))^{1/2}$, a positive invertible
member of $\mathcal D$. If $B$ is invertible, then
\[
 |A_\epsilon B|^2
 =\frac{B^*A^*AB+\epsilon B^*B}{1+\epsilon}.
\]
Apply Lemma~\ref{geo-H-regularization} with $C=B^*A^*AB$ and $D=B^*B$;
its $D$ has finite trace logarithm because it is invertible. The lemma with
$D=1$ also gives $\Delta_\tau(A_\epsilon)\to\Delta_\tau(A)$.
Invertible multiplicativity and these limits prove
$\Delta_\tau(AB)=\Delta_\tau(A)\Delta_\tau(B)$ for arbitrary $A$ and
invertible $B$. Taking adjoints proves the case with the left factor
invertible.

\emph{Zero determinants propagate.} If $\ell(B^*B)=-\infty$ and $A$ is
bounded, put $C=B^*A^*AB$, $D=B^*B$, and $c=\max(1,\|A\|^2)$.
Then $C\leq cD$. Let $P=1_{[0,1/(2c)]}(D)$. Its trace is finite since
$D-1\in L^1(\tau)$. The part of $-\log D$ outside $P$ is trace integrable,
so $\tau(P\log(cD+\eta1)P)\to-\infty$ as $\eta\downarrow0$.
Operator monotonicity of the logarithm gives
\[
 \tau(P\log(C+\eta1)P)
 \leq\tau(P\log(cD+\eta1)P).
\]
If $\log^-C$ had finite trace, the left side would have a finite limit by
monotone spectral convergence on the finite-trace corner $P$; its positive
part is bounded there. This contradiction proves $\ell(C)=-\infty$.
Thus $\Delta_\tau(B)=0$ implies $\Delta_\tau(AB)=0$.

\emph{Both factors arbitrary.} If $\Delta_\tau(B)=0$, the preceding
paragraph proves the desired equality. Otherwise $\ell(B^*B)$ is finite.
Use the same $A_\epsilon$ and the same identity for $|A_\epsilon B|^2$.
The one-invertible-factor case applies to $A_\epsilon B$, and
Lemma~\ref{geo-H-regularization} applies with $D=B^*B$ even though $B$
need not be invertible. Taking limits proves multiplicativity. Finally
$B^*HB=(AB)^*(AB)$, and two applications of multiplicativity, together
with adjoint symmetry, prove the stated answer to the manuscript's question.
\end{proof}

\begin{example}[Refuted: extension to all scalar-plus-compact operators]
\label{geo-H-all-scalars}
On $\ell^2$, the positive operators $2\,1$ and $\tfrac12\,1$ have trace-log
determinants $+\infty$ and $0$. Their product has determinant $1$.
There is no multiplication law with ordinary extended nonnegative numbers
that assigns a definite value to $0\cdot\infty$ and makes all such examples
consistent: replacing the second operator by $\tfrac14\,1$ makes the
product determinant $0$, while the same two determinant values occur.
Replacing it by $\tfrac34\,1$ makes the product determinant $+\infty$.
Moreover $2\,1-2P$ has an infinite positive logarithmic integral and an
infinite negative one, so its unregularized trace logarithm is undefined.
Thus the all-$\lambda$ demand in H.2, with no domain qualification, is false
or undefined, rather than a missing extension of
Theorem~\ref{geo-H-determinant}.
\end{example}

\begin{example}[Refuted: two proposed limiting shortcuts]
\label{geo-H-continuity}
In the diffuse finite algebra $L^\infty([0,1])$ with integral trace, let
$a_n=e^{-n}$ on $[0,1/n]$ and $a_n=1$ elsewhere. Then
$\|a_n-1\|_1=(1-e^{-n})/n\to0$, but
$\Delta(a_n)=e^{-1}\ne\Delta(1)$. Trace-norm continuity, proposed in the
supplied study as an extension argument, is false, even for invertible
approximants. The upper semicontinuity and special regularization actually
proved above are what is needed.

Also $B=-\operatorname{diag}(1/(2n))$ is compact but not trace class,
$1+B>0$, and $\sum_n\log(1-1/(2n))=-\infty$. Its spectral determinant is
$0$, which is finite. Hence ``finite determinant if and only if the compact
perturbation is trace class'' is false when finite includes zero. Finally,
if $B=UL$ is a polar decomposition, then
$B^*HB=L U^*HU L$, not $LHL$. Dropping the conjugating $U$ would assume
part of the multiplicativity being proved. These distinctions are required
by sheets 28--30.
\end{example}

\subsection{A complete reconstruction from the printed 1956 foundations}
\begin{proposition}[Proved: the 1956 determinant ingredients imply the full trace formula]
\label{geo-H-1956-reconstruction}
Grothendieck's printed mixed coefficient estimate, mixed-rank vanishing
rule and determinant zero theorem suffice to deduce
\[
 \operatorname{Tr}T=\sum_j\lambda_j(T)
\]
for every trace-class operator on a complex Hilbert space. The series is
absolutely convergent. The following proof supplies the finite-rank-tail,
zero-free logarithm and global spectral-subspace steps; it is a
reconstruction from the 1956 results, not a claim that these steps are
displayed in that paper.
\end{proposition}
\begin{proof}
Use the symmetric multilinear forms $a_n(T_1,\ldots,T_n)$ normalized by
\[
 D_T(z)=\det(1+zT)=\sum_{n\geq0}z^n a_n(T,\ldots,T),
 \qquad a_1(T)=\operatorname{Tr}T.
\]
The exact printed normalization is \cite[Chapter I, \S2, p.~328,
equations (1)--(4)]{Grothendieck1956Fredholm}: on rank-one arguments the
coefficient is $1/n!$ times the determinant of the pairings.
Chapter~I, \S1, Proposition~1 and its corollary, p.~327, explicitly say
that the mixed exterior product vanishes if $k$ arguments equal an
operator of rank less than $k$. Chapter~II, \S2, Proposition~1,
pp.~346--348, extends the forms continuously to nuclear operators, with
\[
 |a_n(T_1,\ldots,T_n)|\leq\frac1{n!}
       \prod_{i=1}^n\|T_i\|_1
\]
in the Hilbert case: the constant $k_n$ equals $1$ on p.~348. The
trace-class and nuclear norms agree by the singular-value expansion.
Continuity extends the mixed-rank vanishing rule when the other arguments
are arbitrary trace-class operators.

\emph{Minimal exponential type.}
Given $\epsilon>0$, choose a finite-rank operator $A$ with rank $m$ and
write $T=A+R$, where $\|R\|_1<\epsilon/2$. By multilinearity, all terms
with more than $m$ copies of $A$ vanish. For $r=|z|$,
\[
\begin{aligned}
 |D_T(z)|
 &\leq\sum_{n\geq0}\sum_{j=0}^{\min(m,n)}
   \binom nj\frac{r^n}{n!}\|A\|_1^j\|R\|_1^{n-j} \\
 &=e^{r\|R\|_1}\sum_{j=0}^{m}\frac{(r\|A\|_1)^j}{j!}
 \leq C_\epsilon e^{\epsilon r}.
\end{aligned}
\]
The last step absorbs a fixed polynomial into $e^{\epsilon r/2}$.
This obtains minimal exponential type using the printed mixed bound,
without assuming a spectral trace formula or spectral continuity.

\emph{The quasinilpotent case.}
If $\sigma(T)\subseteq\{0\}$, the determinant zero theorem
\cite[Chapter II, \S4, Theorem 2 and corollary, pp.~354--355]{Grothendieck1956Fredholm}
implies that $D_T$ is zero-free. It has an entire logarithm $h$, with
$h(0)=0$ and
\[
 \Re h(z)\leq\epsilon|z|+\log C_\epsilon.
\]
Fix $w$ and a disk of radius $L$ around $w$. Put
$B=\epsilon(L+|w|)+\log C_\epsilon$.
The elementary derivative bound for an analytic function with real part
at most $B$ on this disk gives
\[
 |h'(w)|\leq\frac{2}{L}\bigl(B-\Re h(w)\bigr).
\]
For completeness, when $B-\Re h(w)>0$ apply Schwarz's lemma to
\[
 \frac{(B-h(z))-(B-h(w))}
      {(B-h(z))+\overline{B-h(w)}};
\]
it maps the disk to the unit disk and vanishes at $w$. When the real-part
difference is zero, the maximum principle makes $h$ constant. Let
$L\to\infty$ and then $\epsilon\downarrow0$. Thus $h'=0$, so $D_T=1$,
and its linear coefficient gives $\operatorname{Tr}T=0$.

\emph{The general case.}
First reduce to the separable reducing space
$\overline{\operatorname{ran}T+\operatorname{ran}T^*}$, on whose
orthogonal complement $T$ is zero. Let $M$ be the closed span of the
finite-dimensional primary spaces for all nonzero eigenvalues of $T$.
Choose Jordan flags in these spaces, list them successively, and apply
Gram--Schmidt along the increasing invariant finite-dimensional sums.
This gives an orthonormal basis of $M$ in which $T|_M$ is triangular,
with diagonal the nonzero eigenvalues counted algebraically. Since the
restriction is trace class, its diagonal is absolutely summable, and
\[
 \operatorname{Tr}(T|_M)=\sum_j\lambda_j(T).
\]
No completeness of the generalized eigenvectors in the whole Hilbert
space has been assumed.

The quotient operator on $H/M$, or equivalently the compression $Q$ to
$M^\perp$, is compact and has no nonzero spectral value. To verify this,
if $\lambda\ne0$ lies outside $\sigma(T)$, the bounded inverse of
$T-\lambda$ preserves each primary space and hence $M$, and induces
an inverse on the quotient. If $\lambda\in\sigma(T)\setminus\{0\}$,
the spectral projection of
\cite[Chapter II, \S4, Theorem 4, pp.~359--361]{Grothendieck1956Fredholm}
splits off its finite primary space $E_\lambda\subseteq M$. On the
complementary invariant space $N$ the operator $T-\lambda$ is invertible.
Moreover $M=E_\lambda\oplus(M\cap N)$, and that inverse preserves the
closed span $M\cap N$ of the remaining primary spaces. It therefore
induces an inverse on $N/(M\cap N)\simeq H/M$. This proves
$\sigma(Q)\subseteq\{0\}$.

The compression $Q$ is trace class, so the preceding case gives
$\operatorname{Tr}Q=0$. Taking an orthonormal basis of $M$ followed by
one of $M^\perp$, absolute convergence of trace-class diagonals gives
\[
 \operatorname{Tr}T=\operatorname{Tr}(T|_M)+\operatorname{Tr}Q
                  =\sum_j\lambda_j(T).
\]
This completes the deduction from the printed determinant foundations.
\end{proof}

For comparison, Simon's primary paper explicitly prints the minimal-type
argument in Theorem~3.4, p.~254, the genus-zero entire-function lemma in
Theorem~4.1, pp.~257--258, the product formula in Theorem~4.2 and the trace
identity in Corollary~4.3, p.~258 \cite{catSimon1977}. These are exact
printed locators for his full argument. His use of Grothendieck's
determinant algebra does not itself identify a printed 1956 completion.

\subsection{The historical lead: what the printed sources contain}
\begin{proposition}[Known: published occurrence and the inspected printed treatment]
\label{geo-H-history}
A 1951 \emph{Comptes rendus} note already announces genus zero for
Fredholm determinants of trace operators on Hilbert space, while
promising a detailed treatment elsewhere.
The inspected fourth printing (1966) of the memoir copyrighted 1955
contains two statements independently sufficient for the Hilbert-space
trace formula: genus zero and continuity of the unordered eigenvalue
sequence in actual $\ell^1$. The latter is stated as a Hilbert-space input
to a Banach-space reduction, not proved there. More explicitly,
the published Bourbaki expos\'e~91 states the Hilbert-space genus-zero
product and extends its coefficient identities to trace-class operators;
the text ends with the date June 1957. The 1956
erratum explains that planned Hilbert-space eigenvalue results were removed
from the Fredholm paper, leaving some memoir references inadequate.
The original 1959 Lidskii endnote expressly points to the 1955 memoir
for an indication of the genus-zero claim that it attributes to Weyl.
A later paper, published in 1961 and received on 20 June 1960,
again uses the full Hilbert trace identity as an input.
The audited passages in the memoir and metric \emph{R\'esum\'e} do not
provide the missing full Hilbert proof. The complete 1956 Fredholm article
prints sufficient foundations for Proposition~\ref{geo-H-1956-reconstruction},
but does not display its final minimal-type and global trace deductions.
These facts do not date the crossed-out cote~1 sheets or collate the
memoir's first printing.
\end{proposition}
\begin{proof}[Opened primary texts and exact implication]
The mathematics on the crossed-out sheets, with the repairs just itemized,
proves the trace formula. This does not date the sheets. The opened 1956
paper provides the Fredholm zero and multiplicity theorem used above. Its
introductory note dates an initial draft of that \emph{paper} to 1952 and
its earlier submission to 1953, not cote 1/079--080 or 111--112.
The archive's own catalogue \cite{geoMontpellierCatalogue}, entry for
cote~1, distinguishes undated handwritten notes from a letter dated 1953,
although the dossier is indexed under 1953. Thus that catalogue date
cannot by itself date the individual mathematical sheets.

The earlier primary announcement is explicit
\cite[p.~1557, \S1, final paragraph]{algGrothendieckCR1951Tensor}.
After introducing trace operators and their Fredholm determinants, the
1951 note asserts genus at most one in general and genus zero when the
space is Hilbert, without imposing normality. The same paragraph promises
a detailed Fredholm exposition elsewhere; the three-page note does not
give the Hilbert genus-zero proof. Its first page records presentation on
3 December 1951. This establishes an earlier printed announcement of the
result, without identifying or dating the particular cote~1 proof sheets.

The primary memoir has now been inspected directly
\cite[Chapter II, \S1, no.~4, printed p.~13]{Grothendieck1955Memoir}.
The retained scan is the fourth printing, 1966, with a copyright page dated
1955. On scan page 214, printed p.~II-13, it displays
$\det(1-zu)=1-z\operatorname{Tr}u+\cdots$, states the eigenvalue-zero
correspondence with multiplicities, and explicitly states genus zero for
Hilbert-space operators in the Schatten class with $p\leq1$. Combining
those printed statements yields the product and trace identities by
coefficient comparison. This is a verified printed mathematical occurrence;
it is stronger evidence than the former secondary retrieval lead.

There is a second sufficient assertion in the same memoir,
\cite[Chapter II, \S1, no.~4, Remark 4, pp.~20--22, especially p.~21]{Grothendieck1955Memoir}.
The notation must be read exactly: $\Sigma_{(p)}$ is the metric space of
unordered sequences with summable $p$th powers, whereas $\Sigma_{[p]}$
is the intersection of $\Sigma_{(q)}$ for $q>p$. The Hilbert assertion has
target $\Sigma_{(p)}$ for $0<p\leq1$, hence includes actual unordered
$\ell^1$ continuity at $p=1$. This alone would prove the trace formula:
choose finite-rank $T_n\to T$ in trace norm, write $\lambda(T)$ for the
unordered nonzero eigenvalue sequence padded with zeros, and use
\[
 \left|\sum_j\lambda_j(T_n)-\sum_j\lambda_j(T)\right|
 \leq d_1\bigl(\lambda(T_n),\lambda(T)\bigr)\longrightarrow0.
\]
The inequality holds for every pair of enumerations and then for their
infimum. The asserted target $\Sigma_{(1)}$ supplies absolute convergence.
Finite-dimensional trace gives $\sum_j\lambda_j(T_n)=\operatorname{Tr}T_n$,
and continuity of the trace gives the desired limit. This paragraph is
our explicit deduction from the printed assertion. Remark~4 itself
reduces Banach spectral continuity to this stronger Hilbert assertion;
it supplies no proof of that Hilbert assertion. Thus the memoir contains
another sufficient statement, without thereby containing a second
complete proof of the Hilbert trace formula.

A separate primary text makes the formula explicit:
\cite[\S4(c), printed p.~383, formula (24)]{GrothendieckBourbaki91} states
\[
 \det(1+zu)=\prod_i(1+z\lambda_i)
\]
for Fredholm operators on Hilbert space, with the nonzero eigenvalues
repeated according to their multiplicities and an absolutely summable
sequence. Here ``Fredholm operator'' is defined in \S2, p.~380 as an
operator represented by the projective tensor product, hence is a
trace-class operator in the Hilbert case. The paragraph explicitly extends
formulas (10), p.~378, to this case. Taking their first coefficient gives
$\alpha_1(u)=\sum_i\lambda_i$, and $\alpha_1(u)=\operatorname{Tr}u$
by the definition in \S1. Thus the trace identity is included as the
first of the expressly asserted coefficient identities.
The opening page dates the lecture to March 1954; the final page~384
bears the textual date June 1957. We retain that distinction: the
inspected text documents a June-1957-dated assertion, without inferring
that every sentence was already delivered in March 1954. It announces
the genus-zero result, rather than printing the three-sector proof.

The earlier summary \cite[Chapter II, no.~6, Theorem 8 and Corollary 2,
pp.~107--108]{Grothendieck1952Summary} explicitly gives the trace identity
at the $2/3$-summability threshold in its general locally convex setting.
That is a different statement from the Hilbert-space $p\leq1$ passage
just inspected; no change of threshold is inferred merely from the word
``nuclear''.

The whole accessible memoir was searched after OCR of all 382 scan pages,
and its relevant Chapter~II, \S1, pp.~13--22 were read through. The
genus-zero Hilbert assertion on p.~13 is an input: the corollary on p.~14
explicitly invokes the recalled result; Proposition~1, p.~15, reduces
compositions to that Hilbert assertion; Theorem~4 and its corollaries,
pp.~16--19, use the same reduction to prove the general $2/3$ theorem.
The canonical product on p.~18 retains the linear exponential determined
by the nuclear trace, and does not remove that factor in the full
Hilbert trace-class case. The question on p.~20 concerns general Banach
operators. The other Hadamard occurrences concern sequence-space growth,
and the appended earlier summary repeats the $2/3$ result and asserts
Banach spectral continuity. Remark~4's dependency chain was checked:
Theorem~1, pp.~5--6, concerns common tensor representations of convergent
sequences; Proposition~1, pp.~14--16, supplies the reduction through
Hilbert space. Neither proves Hilbert $\ell^1$ spectral continuity.
The later nuclear-space statement in Chapter~II, \S2, no.~4, p.~65
uses Remark~4 and does not fill this gap. The Appendix to Chapter~I,
pp.~194--195, expressly explains that intended Hilbert eigenvalue results
were removed after the author considered them essentially contained in
Weyl's paper, leaving the earlier references inadequate. No alternative
proof of the Hilbert genus-zero assertion was located in this complete
text audit. In particular the three-sector proof of sheets~111--112 is
not printed in the inspected memoir.

The full 1956 Fredholm article was separately extracted and audited.
Chapter~I, \S3, p.~331, formula~(6) expresses the determinant coefficients
in eigenvalues for finite-rank operators. Chapter~II, \S2,
Proposition~1, pp.~346--348 extends the continuous coefficient and
operator identities; p.~348 notes the better coefficient estimate in
Hilbert space. This extension does not justify passing to an infinite
sum of eigenvalues. Chapter~II, \S4, pp.~354--364 establishes the
zero correspondence, algebraic multiplicities, and resolvent
decomposition. The Hilbert-space remark on pp.~363--364 concerns only
the vanishing of nilpotent parts for normal operators. Chapter~II,
\S5 extends these Fredholm constructions to locally convex spaces;
Chapter~III treats integral kernels. Its final p.~383 recalls order
zero for operators on nuclear spaces, citing the memoir, a different
setting from arbitrary trace-class Hilbert operators. The article does
not carry out the minimal-type argument, the zero-free logarithm argument
or the final global trace deduction. Its mixed coefficient estimate and
the mixed-rank vanishing corollary nevertheless suffice for all these
steps, as Proposition~\ref{geo-H-1956-reconstruction} proves. Thus a
complete proof is reconstructible from the 1956 foundations even though
that completed Hilbert argument is not displayed in the article. The
absence of the manuscript's particular three-sector proof would by itself
have been too narrow a test of this question.

The memoir's initial Hilbert paragraph cites Weyl's 1949 paper.
All four pages of that original paper
have also been opened and visually read \cite{Weyl1949}: they prove
singular-value/eigenvalue majorization, including the infinite-dimensional
extension on p.~411, but do not state or prove the genus-zero Fredholm
conclusion. Consequently one must not turn the memoir's citation into a
claim that Weyl supplied the missing genus proof. Nor does the memoir page
print the trace identity as a separately stated theorem.

The original Russian Lidskii article supplies an independent contemporary
check \cite[p.~487, final discussion, footnote and references 11--12]{algLidskii1959}.
It attributes genus zero under analogous operator hypotheses to Weyl,
using another method based on the results of reference~11, the 1949
paper just inspected. The attached footnote says that an indication
occurs in reference~12, explicitly Grothendieck's \emph{Memoirs of the
American Mathematical Society} no.~16 (1955). Lidskii says that he learned
these facts after preparing the paper for press. The page records receipt
on 23~November~1958, but that does not date the later endnote itself.
This establishes a contemporary pointer to the 1955 memoir, independently
of the inspected 1966 printing. It neither says that Grothendieck printed
a proof nor establishes the correctness of the attribution to Weyl.

Grothendieck's \emph{The trace of certain operators} gives a distinct,
later dated use \cite[pp.~141--143]{algGrothendieckTrace1961}.
The proof of Proposition~(2), remark~(a), p.~142, assumes both absolute
summability of the eigenvalues of a nuclear Hilbert-space operator and
their sum's equality to its trace. Hilbert-space factorization then
answers the cyclic triple-trace question from Sikorski's letter,
introduced on p.~141. The article supplies neither a genus-zero proof
nor an identification or date for an older manuscript. Its receipt line
is 20 June 1960; its reference to Cartier's February 1961 Bourbaki lecture
is expressly marked as added during proof correction. Thus the
1960 receipt and 1961 publication document this later work, without
dating every printed addition or any cote~1 leaf. This continuing
functional-analytic work also prevents assigning the undated sheets
a terminus merely from a broad division of Grothendieck's mathematical
career into periods.

The distinct metric \emph{R\'esum\'e} was also checked, using its official
facsimile reprint and a complete modern transcription as a finding aid
\cite[\S3, no.~6, pp.~53--56; final note 8 bis, p.~77]{algMetricResumeFacsimile}.
The original p.~53 defines $\rho_i(u)$ as the eigenvalues of
$(u^*u)^{1/2}$. Theorem~5, p.~54, gives an eigenvalue-summability criterion
for \emph{Hermitian} $u$, and treats arbitrary $u$ by passing to
$(u^*u)^{1/2}$; its displayed sum is the nuclear norm. The following
Hilbert--Schmidt identities and final note~8 bis concern singular values.
The latter cites Weyl's convexity method and explicitly describes the
1954 S\~ao Paulo seminar as unwritten. These passages do not give the
trace-eigenvalue or genus-zero proof for general nonnormal trace-class
operators. The complete transcription's theorem architecture and the
original relevant pages, introduction, contents and final notes were
checked; no determinant/genus-zero argument was located. This does not
claim a line-by-line visual reading of every facsimile page or collation
of the original issue's cover and publication date.

The contemporary erratum supplies a further primary check:
\cite[p.~119]{Grothendieck1956Erratum} says explicitly that the memoir's
Chapter~II, \S1 referred to Hilbert-space eigenvalue results planned for the
Fredholm article, that Grothendieck subsequently removed them because he
considered them essentially contained in Weyl's paper, and that the
references thereby became inadequate. This confirms the bibliographic
problem in the author's own words. It does not enumerate every removed
statement, and we do not infer that Weyl proves more than the inspected
paper actually contains.
\end{proof}

The published occurrence and the dated assertion are therefore established
independently of any dating of the crossed-out sheets. The manuscript
argument above is a complete mathematical proof after the stated repairs.
The 1951 note already announces Hilbert genus zero and promises a fuller
treatment elsewhere. The inspected memoir states both genus zero and the sufficient $\ell^1$
continuity assertion, without proving those Hilbert inputs there. The
1956 Fredholm article supplies sufficient foundations for the complete
reconstruction just given; its final deduction is not printed there.
The metric \emph{R\'esum\'e} supplies singular-value results, and
Bourbaki \S4(c) announces the genus-zero result. Lidskii's original endnote
corroborates the contemporary pointer to the 1955 memoir while attributing
the claim to Weyl, not a proof to Grothendieck. A claim about exact
first-printing wording would still require collation, and a claim about
the sheet dates requires independent archival evidence. No priority claim
is inferred from these mathematical and bibliographic facts.
